% Generated draft: requires mathematical and source review.
\documentclass[11pt]{article}
\usepackage[margin=25mm]{geometry}
\usepackage{fontspec}
\setmainfont{Latin Modern Roman}
\usepackage{amsmath,amssymb,mathtools,unicode-math}
\setmathfont{Latin Modern Math}
\usepackage{enumitem,needspace,xurl,hyperref}
\setlength{\parindent}{0pt}
\setlength{\parskip}{5pt}
\setlength{\emergencystretch}{3em}
\newcommand{\N}{\mathbb{N}}
\newcommand{\Z}{\mathbb{Z}}
\newcommand{\Q}{\mathbb{Q}}
\newcommand{\R}{\mathbb{R}}
\newcommand{\C}{\mathbb{C}}
\newcommand{\F}{\mathbb{F}}
\newcommand{\A}{\mathbb{A}}
\newcommand{\PP}{\mathbb{P}}
\newcommand{\E}{\mathbb{E}}
\newcommand{\eps}{\varepsilon}
\newcommand{\dd}{\,\mathrm{d}}
\begin{document}
\section*{UnsolvedMath ID: 30005981}
\textbf{Source problem code:} catalogue entry
\subsection{Definitions and mathematical statement}
Let \((X,d)\) be a compact metric space. The supplied record proposes approximating \(X\) by finite graphs \(G_n=(T_n,E_n)\), indexed by integers \(n\geq 0\), where \(T_n\) is the vertex set and \(E_n\) is the edge set. It does not define what it means for these graphs to approximate \((X,d)\).

For a positive real parameter \(p\) and a function \(u:T_n\to\R\), the discrete expression displayed in the source is
\[
\mathcal E_{p,n}(u)=\frac12\sum_{(x,y)\in E_n}|u(x)-u(y)|^p.
\]
Here edges are written as pairs to retain the source formula. Whether each undirected edge occurs once or twice is unspecified; the factor \(1/2\) is retained exactly. Positivity of \(p\) makes this auxiliary expression defined for every real-valued vertex function, but the intended admissible range of \(p\) in the research question is not supplied.

\textbf{Research target, subject to unresolved definitions.} Characterize the compact metric spaces \((X,d)\) that admit a proper scaling constant \(\sigma_p\) for graph approximations whose renormalized discrete \(p\)-energies converge to a meaningful continuum energy. Find further examples beyond the known \(p\)-conductively homogeneous spaces.

This target is not yet a fully specified mathematical predicate. The record does not define a proper scaling constant, its allowed values, or the renormalization sequence determined by it. It also does not specify how a function on \(X\) determines functions on \(T_n\), the continuum function space and energy domain, the meaning of convergence, or the properties required of a meaningful continuum energy. Nor does it settle whether the graph sequence is prescribed, may be chosen, or must satisfy an invariance requirement across choices. These choices must be supplied before a characterization with precise quantifiers can be stated.

The original excerpt motivates the question by referring to the derivative \(\nabla f\), Sobolev functions on Euclidean space, and a proposed characterization of \(W^{1,p}([0,1])\) by convergence of renormalized discrete energies of continuous functions. That passage is textually disordered: its scaling factor appears in the supplied transcription as \texttt{(2p−1 )n}, its discrete energies are not defined there for continuum functions, and the symbol \(n\) is used both in Euclidean-space notation and as a limiting index. Consequently, neither the normalization nor the asserted Sobolev equivalence is adopted here as a precise theorem. The range of \(p\), choice of continuous representatives, finiteness of the limit, and boundary conditions are also unstated.

The supplied dated assessment describes \(p\)-conductive homogeneity as a sufficient structural criterion and mentions Sierpi\'nski-carpet, nested-fractal, and other examples, while explicitly stating that this does not classify all metric spaces. The record does not give the definition or hypotheses of that criterion, or identify the particular examples. These are supplied literature qualifications, not independently verified assertions in this draft.
\subsection{Short English statement}
The problem asks which compact metric spaces support a scaling of graph \(p\)-energies that converges to a continuum energy, and seeks further examples beyond known \(p\)-conductively homogeneous spaces. The source leaves the admissible exponents, graph approximation, scaling rule, convergence, and continuum energy requirements unspecified, so a fully precise characterization problem requires those definitions. Its partially solved label is supplied dated curation.
\subsection{Sources}
\begin{itemize}
\item Supplied local UnsolvedMath registry record, permanent ID \(30005981\), title \emph{Scaling Constants for Discrete Energies on Metric Spaces}, exact source problem code \texttt{catalogue entry}, in the Oberwolfach Reports Open Problems collection. No individual proposer is supplied; the registry records proposed year \(2024\).
\item Supplied report citation: \emph{Recent Developments in Dirichlet Form Theory and Related Fields (2025)}, Oberwolfach Reports, pp.~2425--2510, DOI \texttt{10.4171/owr/2024/42}. \url{https://doi.org/10.4171/owr/2024/42}. The background describes these reports as documenting mathematical workshops and their open-problem discussions. No problem-specific page or speaker attribution is supplied.
\item Supplied literature source: Kigami, \emph{Conductive homogeneity of compact metric spaces and construction of p-energy}. \url{https://ems.press/content/book-files/28253}.
\item Supplied literature source: Cao--Chen, \emph{Whether p-conductive homogeneity holds depends on p}. \url{https://doi.org/10.4171/JFG/156}.
\item Literature assessment date: \texttt{2026-08-22}. The supplied curation labels the problem partially solved, describes \(p\)-conductive homogeneity as a sufficient structural criterion, mentions Sierpi\'nski-carpet, nested-fractal, and other examples, and says that a classification of all metric spaces is not provided. This is supplied dated curation derived from the Oberwolfach material, not an independently verified literature assessment. The source itself requests independent verification of time-sensitive status and cited claims.
\end{itemize}
\end{document}
